\documentclass[11pt]{article}
\usepackage{fontspec}
\usepackage{xeCJK}
\setCJKmainfont[ItalicFont=STSongti-SC-Bold]{Songti SC}
\usepackage{xurl}
\usepackage[a4paper,margin=25mm]{geometry}
\usepackage{amsmath,amssymb,amsthm,mathtools}
\usepackage{xcolor}
\usepackage{booktabs,array,enumitem,fancyhdr,hyperref}
\hypersetup{colorlinks=true,linkcolor=blue!50!black,urlcolor=blue!50!black}
\setlength{\emergencystretch}{2em}
\setlength{\parskip}{0.3em}
\linespread{1.16}
\allowdisplaybreaks
\newtheorem{theorem}{定理}[section]
\newtheorem{proposition}[theorem]{命题}
\newtheorem{lemma}[theorem]{引理}
\theoremstyle{definition}
\newtheorem{definition}[theorem]{定义}
\theoremstyle{remark}
\newtheorem{remark}[theorem]{说明}
\newcommand{\dd}{\,\mathrm d}
\newcommand{\ii}{\mathrm i}
\newcommand{\Rea}{\operatorname{Re}}
\newcommand{\Ima}{\operatorname{Im}}
\newcommand{\Li}{\operatorname{Li}}
\newcommand{\supp}{\operatorname{supp}}
\pagestyle{fancy}
\fancyhf{}
\fancyhead[L]{\small 学习笔记}
\fancyhead[R]{\small 初稿}
\fancyfoot[C]{\thepage}
\setlength{\headheight}{15pt}
\author{}
\date{2026年10月22日初稿}
\makeatletter
\renewcommand{\maketitle}{%
  \thispagestyle{plain}%
  \begin{center}
    {\LARGE\bfseries \@title\par}
    \vspace{0.6em}
    {\small \@date\par}
  \end{center}
  \vspace{0.6em}}
\makeatother
\newcommand{\draftnotice}{\noindent{\small
本文为供逐篇审阅的学习笔记初稿，不是新增研究成果。
文中区分直接证明与引用结果；正文不依赖本文。
本轮已作公式、条件及排版初步核查，尚待作者逐句审定。}\par\medskip}

\XeTeXlinebreaklocale "zh"
\renewcommand{\proofname}{证明}
\date{2026年10月24日初稿}
\title{相信黎曼猜想的理由及其分量}
\begin{document}
\maketitle

\noindent{\small 本篇逐类检查支持 RH 的理由，说明每类理由确切证明了什么、在何处止步。
所引定理注明出处；标为“启发性”的论证不作为证明使用。}\par\medskip

\section{问题}
支持 RH 的数学依据通常可以分为四类：
\begin{enumerate}
 \item[(A)] 已证明的部分结果；
 \item[(B)] 类比情形中已证明的对应命题；
 \item[(C)] 概率与随机矩阵模型；
 \item[(D)] 数值验证（学习笔记第 12 篇已专门讨论）。
\end{enumerate}
本篇依次讨论 (A)(B)(C)，再讨论若干反面考虑。
重点不在于罗列结果，而在于说明每类理由的逻辑分量。
下列部分结果选作比较依据，不是截至本文日期的研究纪录总表；
当代研究路线的系统介绍仍以本书第 22 章为主。

\section{部分结果}
\subsection{临界线上的零点}
记 $N(T)$ 为 $0<\gamma\le T$ 的零点个数，$N_0(T)$ 为其中位于临界线上的个数。
\begin{itemize}
 \item Hardy（1914）：$N_0(T)\to\infty$。
 \item Selberg（1942）：$\liminf N_0(T)/N(T)>0$。
 \item Levinson（1974）：比例大于 $\tfrac13$；Conrey（1989）：大于 $\tfrac25$；
 Pratt、Robles、Zaharescu、Zeindler（2020）：大于 $\tfrac5{12}$。
\end{itemize}
这些结果的分量有一个明确的上限：比例下界即使达到 $1$，也只说明离线零点占零密度，
不能排除任何一个离线零点。

\subsection{零点密度}
记 $N(\sigma,T)$ 为 $\beta\ge\sigma$、$0<\gamma\le T$ 的零点个数。
\begin{itemize}
 \item Bohr--Landau（1914）：对每个 $\delta>0$，$N(\tfrac12+\delta,T)=O_\delta(T)$。
 由于 $N(T)\sim\tfrac{T}{2\pi}\log T$，几乎所有零点都位于临界线的任意窄邻域内。
 \item Ingham（1940）：在 $\tfrac12\le\sigma\le1$ 上，
 $N(\sigma,T)\ll T^{3(1-\sigma)/(2-\sigma)}\log^5T$。
\end{itemize}
这类结果把可能的离线零点限制为稀少的例外，但同样不能排除例外。

\subsection{无零区域}
经典的 Vinogradov--Korobov 型无零区域为（$|t|$ 充分大）
$\sigma\ge1-c(\log t)^{-2/3}(\log\log t)^{-1/3}$。
其宽度随 $t$ 趋于零。现有方法甚至不能证明存在 $\delta>0$，使 $\zeta$ 在 $\Rea s>1-\delta$ 内无零点。

因此，(A) 类结果给出临界线上正比例的零点，或临界线任意固定窄邻域中的密度结论，
却尚未把这些结论加强到每一个零点都在线上。
它们说明零点的统计分布与 RH 相容，但不提供排除个别反例的机制。

\section{类比情形}
\subsection{有限域上的曲线}
设 $C$ 为有限域 $\mathbb F_q$ 上亏格 $g$ 的光滑、射影、几何连通的曲线，$N_m=\#C(\mathbb F_{q^m})$。其 zeta 函数
\[
 Z(C,u)=\exp\Bigl(\sum_{m\ge1}\frac{N_m}{m}u^m\Bigr)=\frac{P(u)}{(1-u)(1-qu)},
\]
$P$ 是 $2g$ 次整系数多项式（Weil 有理性）。Weil（1948）证明 $P$ 的根的模都是 $q^{-1/2}$；
令 $u=q^{-s}$，这正是全部零点满足 $\Rea s=\tfrac12$。Deligne（1974）将其推广到高维簇。
（参考习题第五部分 22.4 计算了一个 $g=1$ 的例子。）

这一类比的分量在于：它表明“Euler 乘积 + 函数方程 + 零点在中线上”的组合在一个完整的平行理论中是真的。
它的局限在于证明的机制：Weil 与 Deligne 的证明依赖 Frobenius 映射及其在上同调上的作用，
零点是一个有限维线性算子的特征值的倒数。对 $\mathbb Q$ 而言，尚无已知的对应算子或几何对象。

\subsection{Selberg zeta 函数}
设 $X=\Gamma\backslash\mathbb H$ 为紧双曲曲面，$0=\lambda_0<\lambda_1\le\lambda_2\le\cdots$ 为 Laplace 算子的特征值。
其 Selberg zeta 函数 $Z_X(s)$ 由闭测地线长度构成 Euler 型乘积。
$Z_X$ 在 $0<\Rea s<1$ 中的零点恰为 $s_j,1-s_j$，其中 $\lambda_j=s_j(1-s_j)$（$j\ge1$），另有 $s=1$ 对应 $\lambda_0$。
于是：
\begin{itemize}
 \item 若 $\lambda_j\ge\tfrac14$，则 $s_j=\tfrac12\pm\ii\sqrt{\lambda_j-\tfrac14}$ 位于中线上；
 \item 若 $0<\lambda_j<\tfrac14$，则 $s_j$ 是 $(0,1)$ 中不等于 $\tfrac12$ 的实数。
\end{itemize}
因此 $Z_X$ 的“RH”成立当且仅当 $X$ 没有 $(0,\tfrac14)$ 中的小特征值。
存在具有小特征值的紧双曲曲面（Randol 1974），故此“RH”并非总成立。

这个例子的意义是双重的。一方面，它实现了 Hilbert--Pólya 设想：零点由一个自伴算子的谱决定，
中线位置等价于谱的下界。另一方面，它表明谱解释本身并不保证零点在中线上，还需要谱的下界。

\subsection{函数方程不足以保证零点位置}
Davenport 与 Heilbronn（1936）构造了模 $5$ 的两个 Dirichlet $L$ 函数的线性组合 $f(s)$，
它满足与 $L$ 函数同型的函数方程，却在 $\Rea s>1$ 中有无穷多个零点。
这个函数没有相应的 Euler 乘积。该例说明，所列函数方程与对称性本身不足以推出零点位置；
一个有效论证必须使用能够排除此类反例的额外信息。
Euler 乘积体现的算术结构是重要候选，但不能由这个例子断言任何可能的证明都必须显式使用它。

\section{概率与随机矩阵模型}
\subsection{Möbius 函数的随机模型}
\begin{proposition}\label{prop:random}
设 $X_1,X_2,\ldots$ 为独立随机变量，$P(X_n=1)=P(X_n=-1)=\tfrac12$，$S_N=\sum_{n\le N}X_n$。
则对每个 $\varepsilon>0$，几乎必然地 $S_N=O(N^{1/2+\varepsilon})$。
\end{proposition}
\begin{proof}
由 Hoeffding 不等式，$P(|S_N|\ge\lambda)\le2\exp(-\lambda^2/(2N))$。
取 $\lambda=N^{1/2+\varepsilon}$，右端为 $2\exp(-N^{2\varepsilon}/2)$，对 $N$ 求和收敛。
由 Borel--Cantelli 引理，几乎必然只有有限个 $N$ 使 $|S_N|\ge N^{1/2+\varepsilon}$。
先对正有理数 $\varepsilon$ 取可数交，再对任意正实数取更小的正有理数，
便可使这些上界在同一个满概率事件上对一切 $\varepsilon>0$ 成立。
\end{proof}
结合学习笔记第 14 篇命题 3.2，若 $\mu(n)$ 在 $\mu(n)\ne0$ 的位置上“像”独立的 $\pm1$，
则 $M(x)=O(x^{1/2+\varepsilon})$，从而 RH 成立。这是 Denjoy 的启发性论证。

其分量有限：$\mu$ 是确定性函数，“像随机变量”没有精确含义；
任何将其精确化的尝试（例如 Sarnak 的 Möbius 随机性猜想）本身都是未解决的问题。
还应区分平方根量级的幂次与固定常数界。独立对称随机游走本身也有
$S_N/\sqrt N$ 无界的正负波动（重对数律），所以这一模型不预测 Mertens 猜想
$|M(x)|<\sqrt x$。该猜想的否定并不反驳本节的
$N^{1/2+\varepsilon}$ 型上界；困难仍在于不能把随机变量的结论直接移用于确定性的 $\mu(n)$。

\subsection{随机矩阵}
Montgomery（1973）在 RH 下证明：对 Fourier 变换支在 $(-1,1)$ 内的检验函数，
规范化零点间距的对关联函数为 $1-\bigl(\frac{\sin\pi u}{\pi u}\bigr)^2$，
与 Gauss 酉系综（GUE）随机 Hermite 矩阵特征值的对关联函数一致。
Odlyzko 对高处零点的大规模计算显示，最近邻间距分布与 GUE 预测高度吻合。

这一吻合支持 Hilbert--Pólya 设想（零点对应于某个自伴算子的谱），从而间接支持 RH。
但须注意逻辑次序：Montgomery 定理以 RH 为前提，GUE 的特征值本身是实数；
统计吻合说明 RH 与已知数据构成一个自洽的整体，而不是独立地推出 RH。

\section{反面考虑}
\begin{enumerate}
 \item \textbf{de Bruijn--Newman 常数}。RH 等价于 $\Lambda\le0$，Rodgers--Tao（2020）证明 $\Lambda\ge0$。
 因此若 RH 成立，则 $\Lambda=0$：$\Xi$ 函数在热方程的反向流下，零点在 $t=0$ 时恰好全部为实数，
 而对任意 $t<0$ 不再如此。Newman 将此表述为：RH 若成立，也只是“恰好”成立。
 \item \textbf{Lehmer 现象}。存在异常接近的相邻零点对（Lehmer 1956）。在热流形变下，
 这样的零点对是最先离开实轴的候选；它们的存在与 $\Lambda=0$ 相呼应。
 \item \textbf{尺度问题}。例如 $S(t)$ 的中心极限定理以 $\log\log T$ 归一化，
 在已连续验证的高度处该量仍不大（学习笔记第 12 篇）。这不是对所有局部统计的统一尺度，
 也不提供离线反例高度的预测。
 \item \textbf{模型的局限}。随机模型可以对“典型”行为作出成功预测，却不能排除个别例外；
 RH 恰恰是一个关于“每一个零点”的命题。
\end{enumerate}

\section{讨论}
四类理由各自都不能证明 RH，它们的分量来自不同的方面：
\begin{itemize}
 \item (A) 给出在线正比例与靠近临界线的密度结果，两者均不同于 RH；
 \item (B) 说明 RH 型命题在有完整结构的平行理论中为真，并指出可能的证明机制（谱或上同调）；
 \item (C) 说明 RH 与零点和 Möbius 函数的统计行为相容；
 \item (D) 排除了低处的反例。
\end{itemize}
它们来源不同，部分论证又相互依赖；综合考察这些相容现象，是形成信念的依据之一。
另一方面，反面考虑表明：若 RH 成立，它是在一个没有余量的临界位置上成立；
而 Davenport--Heilbronn 的例子提醒我们，函数方程的形式对称性不足以代替更强的结构信息。
对一个关于“每一个零点”的命题而言，所有统计性或类比性的支持都不能替代证明。

\section*{参考}
\begin{itemize}
 \item E.\ C.\ Titchmarsh, \emph{The Theory of the Riemann Zeta-Function}, 2nd ed., 1986, Ch.~IX, X, XIV.
 \item K.\ Pratt, N.\ Robles, A.\ Zaharescu, D.\ Zeindler, More than five-twelfths of the zeros of $\zeta$ are on the critical line,
 \emph{Res.\ Math.\ Sci.} 7 (2020). \url{https://arxiv.org/abs/1802.10521}
 \item A.\ Weil, \emph{Sur les courbes algébriques et les variétés qui s'en déduisent}, Hermann, 1948；
 P.\ Deligne, La conjecture de Weil I, \emph{Publ.\ Math.\ IHÉS} 43 (1974).
 \item H.\ Davenport, H.\ Heilbronn, On the zeros of certain Dirichlet series, \emph{J.\ London Math.\ Soc.} 11 (1936).
 \item H.\ L.\ Montgomery, The pair correlation of zeros of the zeta function, \emph{Proc.\ Sympos.\ Pure Math.} 24 (1973).
 \item A.\ M.\ Odlyzko, On the distribution of spacings between zeros of the zeta function, \emph{Math.\ Comp.} 48 (1987).
 \item B.\ Rodgers, T.\ Tao, The de Bruijn--Newman constant is non-negative, \emph{Forum Math.\ Pi} 8 (2020).
 \item P.\ Sarnak, Three lectures on the Möbius function, randomness and dynamics, IAS 讲义, 2011.
 \item B.\ Randol, Small eigenvalues of the Laplace operator on compact Riemann surfaces, \emph{Bull.\ AMS} 80 (1974).
\end{itemize}
\end{document}
